\section{The component \texorpdfstring{$\PSL_4(5)$}{PSL4(5)}}\label{sec:exceptional}

The square-discriminant six-dimensional space over $\F_5$ has plus
type, and
\[
 \mathrm P\Omega_6^+(5)\cong\PSL_4(5).
\]
This is the case excluded from Proposition~\ref{prop:orthogonal}.
In a mixed frame the norm colors occur in counts two and four.
Its inner sign group has rank three and does not distinguish the two
positions of the smaller color class. Thus it meets neither hypothesis
of Lemma~\ref{lem:orthogonal-centralizer}.
We shall first use any available outer direction to repair this
defect. If no suitable extension is available, an equivalent-poset
construction will retain the homology of the component and its
centralizer.

\subsection{The two natural representations}

Put $L=\PSL_4(5)$ and $V=\F_5^4$. The determinant identifies
$\PGL_4(5)/L$ with $\F_5^\times$, since fourth powers in $\F_5^\times$
are trivial. Let $d$ denote the outer class of
$\diag(2,1,1,1)$, and let $\gamma$ be inverse transpose. The
automorphism theorem for $L$
\citep[Section~2 and statements~3.2--3.6]{Steinberg1960} gives
\begin{equation}\label{eq:psl45-outer}
 \Aut(L)=\PGL_4(5)\rtimes\langle\gamma\rangle,\qquad
 \Out(L)=\langle d,\gamma\mid
 d^4=\gamma^2=1,\ \gamma d\gamma=d^{-1}\rangle.
\end{equation}
There is no field automorphism because the field is prime. Thus the
outer group is dihedral of order eight. We need to identify precisely
the subgroup coming from six-dimensional isometries.

Choose a basis $e_1,\ldots,e_4$ of $V$, write
$e_{ij}=e_i\wedge e_j$, choose a volume form, and put
$W=\bigwedge^2V$. Define the
nondegenerate symmetric form $\beta$ by
\[
 u\wedge v=\beta(u,v)\,\mathrm{vol}\qquad(u,v\in W).
\]
The pairs of basis vectors
\[
 (e_{12},e_{34}),\quad(e_{13},e_{24}),\quad(e_{14},e_{23})
\]
are hyperbolic pairs, with respective pairings $1,-1,1$.
Consequently $W$ has plus type and square discriminant over $\F_5$.
For $g\in\GL(V)$,
\begin{equation}\label{eq:psl45-wedge-multiplier}
 \beta((\bigwedge^2g)u,(\bigwedge^2g)v)=\det(g)\beta(u,v).
\end{equation}
The standard exterior-square covering
$\SL_4(5)\to\Omega_6^+(5)$ has kernel $\{\pm I\}$
\citep[Theorem~12.20 and Corollary~12.21]{Taylor1992}; after
projectivization it is the displayed isomorphism with $L$.
In particular, the two groups called $L$ in these representations
are identified.

Let $\tau$ interchange complementary basis vectors with the exterior
signs:
\[
 e_{12}\leftrightarrow e_{34},\qquad
 e_{13}\leftrightarrow-e_{24},\qquad
 e_{14}\leftrightarrow e_{23}.
\]
It is an isometry, $\tau^2=1$, and $\det(\tau)=-1$. Directly from
the wedge pairing,
\[
 \tau(\bigwedge^2g)\tau^{-1}
       =\det(g)\bigwedge^2(g^{-T}).
\]
Thus $\tau$ induces $\gamma$ on the projective image of $\SL_4(5)$.
The exterior-square image of $\PGL_4(5)$ and $\tau$ together give
the full projective similarity group of $W$
\citep[Theorem~12.18]{Taylor1992}: the former has
relative determinant
$\det(\bigwedge^2g)/\det(g)^3=1$, whereas $\tau$ has relative
determinant $-1$; their order is that of this similarity group.
Equivalently, this is the natural $D_3=A_3$ realization of
\eqref{eq:psl45-outer}. We write $\PO(W)$ for the projective image
of the isometry group, a subgroup of $\Aut(L)$.

A similarity can be rescaled to an isometry exactly when its
multiplier is square. Equation~\eqref{eq:psl45-wedge-multiplier}
and the fact that $\tau$ is an isometry therefore show that
\begin{equation}\label{eq:psl45-natural-outer}
 T:=\PO(W)/L=\langle d^2,\gamma\rangle\cong C_2^2.
\end{equation}
Here $L$ is the joint determinant--spinor-norm kernel in $\PO(W)$.
In this dimension $-I$ has determinant and spinor norm one, so both
characters descend to the projective group. In particular, there is
no ambiguity from changing an isometry representative by $-I$.

There are only two possibilities for an elementary abelian subgroup
$B\leq\Out(L)$ satisfying $B\cap T=1$. Since $T$ has index two,
either $B=1$, or $B$ has order two and is generated by
$d\gamma$ or $d^3\gamma$. These are precisely the graph cosets
whose diagonal determinant class is nonsquare.

The remaining preparatory inputs concern centralizers and homology.
They will distinguish the three branches of the final argument.

\subsection{An inheritance fact for centralizers}

The passage from an extension to one of its outer directions requires
the following elementary observation.

\begin{lemma}\label{lem:psl45-intermediate-core}
Let $J\trianglelefteq H\leq G$ be finite groups, and suppose that
$H/J$ is a $p$-group. If $O_p(C_G(H))=1$, then
$O_p(C_G(J))=1$.
\end{lemma}

\begin{proof}
Put $Q=O_p(C_G(J))$. The group $H$ normalizes $C_G(J)$ and its
characteristic subgroup $Q$. Since $J$ centralizes $Q$, the action
of $H$ on $Q$ factors through the $p$-group $H/J$.
If $Q\ne1$, orbit counting on the underlying set of $Q$ gives
\[
 |C_Q(H)|\equiv |Q|\equiv0\pmod p.
\]
The identity is fixed, so $C_Q(H)\ne1$. Moreover $C_G(H)$ is a
subgroup of $C_G(J)$, normalizes $Q$, and centralizes the action of
$H$. Hence
$C_Q(H)=Q\cap C_G(H)$ is a nontrivial normal $p$-subgroup of
$C_G(H)$, a contradiction.
\end{proof}

\subsection{Homology of the component}

The final branch will use nonzero homology of $L$ itself. We prove
that fact directly, without applying the assertion under induction.

\begin{lemma}\label{lem:psl45-homology}
For $L=\PSL_4(5)$,
\[
 \widetilde\chi(\mathcal A_2(L))\equiv-1\pmod5.
\]
In particular $\widetilde H_*(\mathcal A_2(L);\Q)\ne0$.
\end{lemma}

\begin{proof}
We first show that every elementary abelian $2$-subgroup $E\leq L$
has rank at most four. Apply Lemma~\ref{lem:projective-rank} to
$E\leq L\leq\PGL_4(5)$, with $p=2$ and $n=4$. If $2a$ is the
rank of its scalar-commutator pairing, then $2^a\mid4$ and
\begin{equation}\label{eq:psl45-rank}
 \rk E\leq2a+\frac4{2^a}-1.
\end{equation}
The possibilities $a=0,1,2$ give the bounds $3,3,4$, respectively.

Next let $x\in L$ be an involution and choose a lift
$X\in\SL_4(5)$, with $X^2=\lambda I$.
If $\lambda$ were nonsquare, $T^2-\lambda$ would be irreducible
over $\F_5$. Its two roots would each have multiplicity two in
$X$, giving
\[
 \det X=\lambda^2=-1,
\]
a contradiction. Thus $\lambda=a^2$ for some $a\in\F_5^\times$.
Replacing $X$ by $a^{-1}X$ keeps determinant one, since $a^4=1$,
and makes $X^2=I$. The lift is noncentral, and determinant one
forces its $1$- and $-1$-eigenspaces to have dimensions two and two.
A projective centralizer of $x$ preserves or exchanges these two
spaces: if $YXY^{-1}=cX$, then the eigenvalues force $c=\pm1$.
The $5$-part of this block stabilizer has order at most
\[
 |\GL_2(5)|_5^2=5^2.
\]
Passing to a subgroup and then to its projective quotient cannot
increase this bound. Hence $|C_L(x)|_5\leq5^2$.

A Sylow $5$-subgroup $P$ of $L$ has order $5^6$.
If $P$ normalized a nontrivial elementary abelian $2$-subgroup $E$,
its conjugation image would lie in $\GL_{\rk E}(2)$. By
\eqref{eq:psl45-rank}, the $5$-part of that group has order at most
five. It would follow that $|C_P(E)|\geq5^5$, contradicting the
bound for $C_L(x)$ for any $1\ne x\in E$.

Thus $P$ fixes no vertex of $\mathcal A_2(L)$. A simplex fixed
setwise by $P$ would have each vertex fixed, because an automorphism
of a finite chain fixes its ordered vertices. There are therefore
no fixed nonempty simplices. Every orbit of nonempty simplices
has cardinality divisible by five, so ordinary Euler
characteristic is zero modulo five. Reduced Euler characteristic
is one less, giving the stated congruence.
Euler characteristic is the alternating sum of rational homology
dimensions, so the nonvanishing conclusion follows.
\end{proof}

\subsection{The equivalent poset used in the final branch}

We record the precise specialization of the equivalent-poset
theorem that will be needed. A product below includes the trivial
subgroup in either coordinate, but excludes $(1,1)$.

\begin{theorem}[Equivalent poset for a centerless component]
\label{thm:psl45-equivalent-poset}
Let $G$ be a finite group, let $p$ be a prime, and let $L$ be a
centerless component of $G$ with $p\mid |L|$. Put $K=C_G(L)$ and
\[
 \mathcal Q=
 \bigl((\mathcal A_p(L)\cup\{1\})\times
             (\mathcal A_p(K)\cup\{1\})\bigr)\setminus\{(1,1)\},
 \qquad
 \mathcal F=\{F\in\mathcal A_p(G):F\cap LK=1\}.
\]
Give $\mathcal Q$ the coordinatewise order and $\mathcal F$ the
inclusion order. Define
\[
 \mathcal F_0=
 \{F\in\mathcal F:F\leq N_G(L),\
                         O_p(C_G(LF))\ne1\}.
\]
The disjoint union $\mathcal X=\mathcal Q\sqcup\mathcal F$ is
ordered by these orders and by the following crossed comparisons:
\begin{equation}\label{eq:psl45-cross}
 F<(U,V)
 \quad\Longleftrightarrow\quad
 \begin{cases}
 C_{UV}(F)\ne1,\\
 C_{UV}(F)\not\leq L&\text{if }F\in\mathcal F_0.
 \end{cases}
\end{equation}
There are no comparisons directed from $\mathcal Q$ to
$\mathcal F$. These rules define a poset, and
$\mathcal X\simeq\mathcal A_p(G)$.
\end{theorem}

This centerless case is supplied by
\citet[Propositions~4.11--4.12]{PitermanSmith2025}.
It is also the specialization of
\citet[Definition~3.2 and Theorem~3.3]{Piterman2026}, together with
the equivalence preceding that definition, to $B=1$ and
$\mathcal P_L=\mathcal A_p(L)$.
Its component hypothesis is $p\mid|L|$ and $p\nmid|Z(L)|$;
the centerless hypothesis here supplies the latter.
The required condition $B\cap LC_G(L)=1$ is automatic for $B=1$.
No normality of $L$ in $G$, and no defining-characteristic assumption,
is required. Since $Z(L)=1$, the product $LK$ in the statement is
the direct product $L\times K$, so the products $UV$ used in
\eqref{eq:psl45-cross} are unambiguous.

\subsection{Exclusion of a minimal counterexample}

The preparatory results are now in place. The first two branches
return to frame cycles; the third uses the equivalent poset. We use \emph{configuration} in the sense of
Definition~\ref{def:faithful-configuration}: $A\leq N_G(L)$ is elementary abelian,
$A\cap L\ne1$, the group $H=LA$ acts faithfully on $L$, and
$O_2(C_G(H))=1$.

The first branch starts from a six-dimensional sign configuration
whose rank and distinct weights permit propagation. If none exists,
Lemma~\ref{lem:psl45-intermediate-core} will exclude every faithful
extension with trivial centralizer core whose outer image meets $T$
nontrivially. The only possible remaining nontrivial outer image is
then an odd-diagonal graph line, which gives the second branch.
If that branch is absent as well, the equivalent-poset theorem and
the component homology proved above give the final argument.
These transitions are proved inside the proposition; in particular,
the graph-line restriction is imposed only after the first branch
has been excluded.

\begin{proposition}\label{prop:psl45}
A minimal-order counterexample to the rational homological
implication at $p=2$ has no simple component isomorphic to
$\PSL_4(5)$.
\end{proposition}

\begin{proof}
Suppose that $G$ is such a counterexample with component $L$.
Faithful extensions of $L$ in $N_G(L)$ will be identified with
their images in $\Aut(L)$. This identifies the original copy of
$L$ with the inner automorphism subgroup and allows us to use
either of its two natural representations.

\paragraph{First branch: a six-dimensional sign configuration.}
Let $\mathcal D$ be an orthogonal frame of $W$ with two or four
nonsquare-norm lines; call such a frame mixed. Its projective sign
group
\[
 S(\mathcal D)\cong\F_2^6/\langle(1,1,1,1,1,1)\rangle
\]
has rank five. Write $t$ for total sign parity and $c$ for parity
on the nonsquare-norm positions. Both descend to the projective
sign group, because both color classes have even size.
The determinant and spinor-norm conditions give
\begin{equation}\label{eq:psl45-inner-signs}
 S_L(\mathcal D):=S(\mathcal D)\cap L=\ker t\cap\ker c,
 \qquad \rk S_L(\mathcal D)=3.
\end{equation}
The two characters are independent. Consequently the image of
$S(\mathcal D)$ in $T$ is all of $T$.

Consider configurations $A$ such that, in the six-dimensional
representation,
\[
 S_L(\mathcal D)\leq A\leq S(\mathcal D),\qquad
 \rk A\geq4,
\]
for some mixed frame $\mathcal D$, and no two coordinate weights
are equal on $A$. Suppose first that this family is nonempty,
and choose $A$ of maximal rank in it. Put $H=LA$.

We check the centralizer assertion that makes this a permissible
maximizing family. Choose a linear section of the binary sign
preimage of $A$, and let $w_1,\ldots,w_6$ be its six coordinate
characters. Their differences span $A^*$, because the projective
sign action is faithful, and they are distinct by assumption.
Any projective similarity centralizing $A$ permutes these
characters by one common translation $\chi$: its scalar
commutators with sign lifts are $\pm1$ and form the character
$\chi$. If $\chi\ne0$, the six characters form three pairs
$\{w,w+\chi\}$. The affine span then has dimension at most three,
as it is spanned by two differences of pair representatives and
$\chi$. This contradicts $\rk A\geq4$.

It follows that every such centralizer preserves the six lines.
If its diagonal entries are $a_i$ and its multiplier is $\mu$,
then $a_i^2=\mu$ for every $i$. Rescaling by $a_1^{-1}$ makes
all entries $\pm1$. By the full similarity realization of
$\Aut(L)$ above, we have proved
\begin{equation}\label{eq:psl45-six-centralizer}
 C_{\Aut(L)}(A)=S(\mathcal D).
\end{equation}
In particular maximality gives $A=S(\mathcal D)\cap H$.
Any elementary overgroup of $A$ centralizes $A$ and normalizes $L$:
a nonidentity element of $A\cap L$ cannot also lie in a distinct
conjugate component. If the larger group gives a configuration,
its faithful image is therefore contained in the same
$S(\mathcal D)$ by \eqref{eq:psl45-six-centralizer}; it still
contains $S_L(\mathcal D)$, separates positions, and has rank at
least four. It belongs to our family. Adjoining an involution of
$C_L(A)$ also preserves $H$ and its centralizer and gives a
member of the same family. Maximality thus supplies both
centralizer and strict-overgroup conditions of
Lemma~\ref{lem:propagation}.

It remains to obtain a supported coefficient at another member
of this same maximal family. Let $r=\rk(H/L)$; then $r=1$ or $2$,
and $\rk A=3+r$. For any mixed frame $\mathcal D'$,
\[
 |S(\mathcal D'):S(\mathcal D')\cap H|=2^{2-r},
 \qquad L(S(\mathcal D')\cap H)=H.
\]
Indeed its full signs map onto $T$. For an arbitrary frame, the
image of its signs is a subgroup of $T$, so this index is at most
$2^{2-r}$. Mixed frames consequently attain the maximum index.

If $r=2$, all five sign directions survive and distinguish the
positions. If $r=1$, the intersection imposes one of
$t=0$, $c=0$, or $t+c=0$. The first condition distinguishes all
six positions. Either of the others is parity on one color class,
and distinguishes positions precisely when that class has four
positions, rather than two. To see this directly, equality of
positions $i,j$ on the intersection would mean that the
two-coordinate functional $x_i+x_j$ is its single defining
functional.

Theorem~\ref{thm:frame-bound} supplies a nonzero undecorated sign
cycle using all frames of $W$. Some mixed frame has a nonzero
coefficient. In fact, if all mixed-frame coefficients vanished,
take a binary merger tree at a homogeneous frame and choose a
bottom pair. Over $\F_5$, the plane on two lines of the same norm
type has exactly two unordered orthogonal frames, one of each
homogeneous color. The plane relation expresses the given
coefficient, up to its nonzero orientation sign, as the coefficient
after replacing that pair by its opposite-color frame. The latter
frame is mixed. All homogeneous-frame coefficients would
therefore vanish as well, contradicting the nonzero cycle.

If necessary, we can change the color counts of the selected
mixed frame before restricting the cycle. There is a similarity
of $W$ with nonsquare multiplier: in hyperbolic coordinates,
multiply one member of each hyperbolic pair by $2$ and fix the
other member. It normalizes $\PO(W)$ and $L$, and exchanges the
two norm types on every line. On determinant--spinor coordinates it acts by
$(t,c)\mapsto(t,c+t)$, since the spinor norm of a reflected line is
multiplied by the similarity multiplier. Apply this automorphism to the
entire sign cycle while keeping $H$ fixed. We do not need it to
normalize $H$: every transformed top group still has image contained
in $T=\PO(W)/L$, and a mixed frame still has image all of $T$.
Its intersection index is consequently at most $[T:H/L]$, with
equality at the selected mixed frame. The resulting cycle
has a nonzero coefficient at a mixed frame for which the parity
condition just described, if it is a color condition, uses four
positions.

The maximum-index form of Lemma~\ref{lem:restriction} now applies
to this actual coefficient, with cone rank five and
$c=2-r\leq1$. It produces a cycle in $\mathcal A_2(H)$ with a
nonzero full-flag coefficient at
$A'=S(\mathcal D')\cap H$. This subgroup has rank $3+r=\rk A$,
separates positions, and generates the same $H$ over $L$.
It is therefore in the maximizing family, with the same
centralizer $C_G(H)$, and satisfies both propagation conditions.
Lemma~\ref{lem:propagation} gives the desired contradiction in
this branch.

\paragraph{Passage to the remaining outer directions.}
We may now assume that the family in the first branch is empty.
This has a stronger consequence than merely excluding the
particular sign groups. Suppose $H_1>L$ is a faithful elementary
$2$-extension in $N_G(L)$, has $O_2(C_G(H_1))=1$, and its outer
image is a nontrivial subgroup of $T$. Choose a mixed frame
whose larger color class is the one required by the possible
single parity condition. Its sign intersection with $H_1$ has
rank at least four, separates positions, maps onto $H_1/L$,
and contains $S_L(\mathcal D)$. It would be a configuration in
the excluded family.

More generally, let $H_1$ be any faithful elementary $2$-extension
with $O_2(C_G(H_1))=1$. If its outer image $B$ met $T$
nontrivially, take the intermediate subgroup $J$ above an
order-two subgroup of $B\cap T$. The elementary complement
defining $H_1$ supplies an order-two complement for $J$.
Moreover $J\trianglelefteq H_1$ and $H_1/J$ is a $2$-group,
so Lemma~\ref{lem:psl45-intermediate-core} gives
$O_2(C_G(J))=1$. This contradicts the preceding paragraph.
Thus every such extension has
\begin{equation}\label{eq:psl45-avoid}
 (H_1/L)\cap T=1.
\end{equation}
By \eqref{eq:psl45-outer}--\eqref{eq:psl45-natural-outer}, a
nontrivial outer image is now a single odd-diagonal graph line.

\paragraph{Second branch: an odd-diagonal graph extension.}
Suppose that such an extension $H=L\langle b\rangle$ exists,
where $b$ is an involution, its outer class is $d\gamma$ or
$d^3\gamma$, and $O_2(C_G(H))=1$.

We describe its compatible frames explicitly in $V$. Write
$b=\operatorname{Int}(B)\gamma$ with $B\in\GL_4(5)$.
The condition $b^2=1$ says $BB^{-T}$ is scalar, whence
$B^T=\varepsilon B$ with $\varepsilon=\pm1$.
If $\varepsilon=-1$, the determinant of $B$ is the square of
its Pfaffian. This is impossible in an odd-diagonal graph coset.
Thus $M=B^{-1}$ is a nonsingular symmetric Gram matrix with
nonsquare determinant. The linear centralizer of $b$ consists
projectively of the similitudes of $M$: the commuting equation
is $gBg^T=\lambda B$, equivalently
$g^TMg=\lambda M$. An orthogonal nondegenerate line frame for
$M$ has one or three nonsquare-norm lines, since the product of
its four line norms is nonsquare.

For such a frame $\mathcal D$, let $S=S(\mathcal D)$ be the full
projective coordinate signs in $\PGL_4(5)$, and put
\[
 S_L=S\cap L,\qquad E=S\times\langle b\rangle,\qquad
 A=E\cap H=S_L\times\langle b\rangle.
\]
Here $\rk S=3$, and $S_L$ consists of the even signs and has
rank two. The image of $S$ in $\Out(L)$ is
$\langle d^2\rangle$, whereas $H/L$ is an odd-diagonal graph
line. Therefore $\rk E=4$, $\rk A=3$, $|E:A|=2$, and $LA=H$.
These statements hold for every compatible pair $(\mathcal D,b)$
with $b$ in this fixed extension.

There is a special centralizer check in this dimension. In frame
coordinates, lifts
\[
 \diag(-1,-1,1,1),\qquad \diag(-1,1,-1,1)
\]
generate $S_L$ projectively. Their four joint weights are exactly
the four points of $\F_2^2$. A projective linear centralizer of
$S_L$ acts on these weights by a translation. Centralizing $b$
also makes it a similitude of $M$. A nonsquare multiplier would
exchange the two line colors, impossible because their
cardinalities are one and three. A square multiplier preserves
the unique minority-color line. A nonzero translation of
$\F_2^2$ fixes no point, so the translation must be zero.
The centralizer thus preserves every coordinate line. Its
diagonal entries have a common square by the similitude equation,
so projectively they are signs. We have proved
\begin{equation}\label{eq:psl45-four-centralizer}
 C_{\PGL_4(5)}(A)=S,\qquad C_L(A)=S_L\leq A.
\end{equation}
In particular the order-two centralizer condition for propagation
holds for every pair under consideration.

We next construct the cycle, allowing all involutions $b\in H$
in its fixed nontrivial outer coset and all their compatible
orthogonal frames. For each fixed $b$, these are exactly the
frames of the symmetric four-space just described. The fraction
in Theorem~\ref{thm:frame-bound}, at $u=5$ and $h=4$, is
\[
 x^4+(1-x)^4-\frac{2}{24^2}=\frac14,
 \qquad x(1-x)=\frac5{24}.
\]
This is uniform in the symmetric form and its compatible frames.

For a fixed frame and one compatible $b$, there are at least
sixteen choices of its graph axis within $H$. Indeed let
$t=\diag(t_1,t_2,t_3,t_4)$ in these coordinates, with
$t_i\in\F_5^\times$ and $\prod_i t_i=1$.
Modulo scalar matrices these give
\[
 \frac{4^3}{4}=16
\]
distinct elements of $L$. Since $b$ inverts the diagonal torus,
every $tb$ is an involution; it remains in $H$ and preserves
compatibility with the frame. Its symmetric form is again
diagonal and has nonsquare determinant. The sixteen resulting
order-two subgroups $\langle tb\rangle$ are distinct.

The initial coefficients are detected before imposing the
axis-removal equations. The full group $E$ recovers its linear
kernel $E\cap\PGL_4(5)=S$, hence its coordinate frame. In the
split chains of Lemma~\ref{lem:decorations}, take flags through
the decoration line $\langle b\rangle$ and then grow along the
sign part. The initial line determines the graph axis, and the
remaining flag recovers the original sign coefficient. Different
compatible data therefore retain independent coefficient
parameters. Since
\[
 \frac14>\frac1{16},
\]
Lemma~\ref{lem:decorations} gives a nonzero cycle with cone groups
$E$ of rank four. The intersection index with $H$ is constantly
two. Lemma~\ref{lem:restriction}, with $c=1<4$, produces a cycle
in $\mathcal A_2(H)$ having a nonzero full-flag coefficient at
one of the groups $A=E\cap H$ of rank three.

For completeness, no strict elementary overgroup $D>A$ can
give a configuration. Such a $D$ normalizes $L$, as before.
If $LD$ is faithful and $O_2(C_G(LD))=1$, its outer image
satisfies \eqref{eq:psl45-avoid}. It contains the outer image of
$A$, so the dihedral-group calculation forces it to equal that
same order-two subgroup. Every element of $D$ may therefore
be multiplied by either $1$ or $b$ to lie in $L$. The resulting
element still centralizes $A$, and hence lies in $S_L$ by
\eqref{eq:psl45-four-centralizer}. Thus $D=A$, a contradiction.
The supported subgroup satisfies both propagation conditions,
so Lemma~\ref{lem:propagation} excludes this branch.

\paragraph{Third branch: nontrivial centralizer cores.}
We are left with the assertion that every faithful elementary
$2$-extension $H_1>L$ in $N_G(L)$ satisfies
$O_2(C_G(H_1))\ne1$. Put $K=C_G(L)$. By
Lemma~\ref{lem:component-core}, $O_2(K)=1$. Also $K<G$,
since $L$ is nonabelian. Minimality gives
\begin{equation}\label{eq:psl45-K-homology}
 \widetilde H_*(\mathcal A_2(K);\Q)\ne0,
\end{equation}
with the augmented convention if this poset is empty.

Apply Theorem~\ref{thm:psl45-equivalent-poset} to this centerless
component of even order. In its notation, no point of
$\mathcal F$ is comparable with a pure $L$-point $(U,1)$.
To prove this, suppose that the first condition in
\eqref{eq:psl45-cross} holds for such a pair and choose
$1\ne x\in C_U(F)$. Every element of $F$ fixes $x\in L$ and
therefore normalizes $L$: distinct conjugate simple components
have trivial intersection. Since $F\cap LK=1$, the group
$LF$ acts faithfully on $L$. Indeed an element $lf$ of its
kernel would give $f\in LK$, forcing $f=1$ and then $l=1$.
Moreover $F$ is a nontrivial elementary complement to $L$.
The branch assumption implies $F\in\mathcal F_0$.
But then $C_U(F)\leq L$ violates the second condition of
\eqref{eq:psl45-cross}. This proves the no-comparison assertion.

Let $Y=\Delta(\mathcal Q)$, and let $Z$ be the full subcomplex
of $\Delta(\mathcal X)$ obtained by omitting all pure $L$-points.
A simplex containing such a point contains no vertex of
$\mathcal F$, by the preceding paragraph. Consequently
\[
 \Delta(\mathcal X)=Y\cup Z,\qquad
 Y\cap Z=\Delta(\mathcal Q'),\qquad
 \mathcal Q'=\{(U,V)\in\mathcal Q:V\ne1\}.
\]
If $\mathcal A_2(K)$ is nonempty, the inclusion
$\mathcal Q'\hookrightarrow\mathcal Q$ is nullhomotopic.
For any fixed $U_0\in\mathcal A_2(L)$, the following pointwise
comparisons give the required homotopy to a constant:
\[
 (U,V)\ \geq\ (1,V)\ \leq\ (U_0,V)\ \geq\ (U_0,1).
\]
The reduced Mayer--Vietoris sequence now gives an injection
\begin{equation}\label{eq:psl45-MV}
 \widetilde H_n(Y;\Q)\longrightarrow
                    \widetilde H_n(\Delta(\mathcal X);\Q).
\end{equation}
Indeed the first coordinate of
$\widetilde H_n(Y\cap Z)\to
 \widetilde H_n(Y)\oplus\widetilde H_n(Z)$ is zero, so its image
contains no nonzero element of the form $(y,0)$.

The deleted-bottom product poset $\mathcal Q$ has the homotopy
type of the join
$\mathcal A_2(L)*\mathcal A_2(K)$, by the deleted-bottom product/join
description in Lemma~\ref{lem:product-shuffle}. Over $\Q$ its augmented join
formula is
\[
 \widetilde H_n(Y;\Q)\cong
 \bigoplus_{i+j=n-1}
 \widetilde H_i(\mathcal A_2(L);\Q)\otimes_\Q
 \widetilde H_j(\mathcal A_2(K);\Q).
\]
Lemma~\ref{lem:psl45-homology} and
\eqref{eq:psl45-K-homology} make this nonzero in some degree.
If $\mathcal A_2(K)$ is empty, then
$\mathcal Q=\mathcal A_2(L)$ and $\mathcal Q'=\varnothing$.
There are no comparisons at all between $\mathcal Q$ and
$\mathcal F$, so its nonzero reduced homology injects directly
into that of $\mathcal X$; the same join formula uses the
degree $-1$ class of the empty factor. Thus this case also
gives the conclusion of \eqref{eq:psl45-MV}.

Finally $\mathcal X\simeq\mathcal A_2(G)$, so the nonzero
class contradicts the assumed rational acyclicity of
$\mathcal A_2(G)$. All three branches have been excluded.
\end{proof}
